\documentclass{article}
\usepackage[T1]{fontenc}
\usepackage{lmodern}
\usepackage{graphicx}
\usepackage{amsmath,amssymb}
\usepackage[a4paper,margin=2cm]{geometry}
\usepackage[absolute,overlay]{textpos}
\setlength{\TPHorizModule}{1mm}
\setlength{\TPVertModule}{1mm}
\usepackage[indonesian]{babel}
\usepackage{hyperref}
\setlength{\emergencystretch}{2em}
% unit: Halaman 70

\begin{document}

% === Halaman 70 (python/native) ===

Alasan mengapa huruf  J dibuang:

1. Matriks Playfair hanya 5×5

Playfair cipher menggunakan tabel 55, jadi hanya bisa memuat 25 huruf.  Sementara alfabet latin 

ada 26 huruf (A–Z). Artinya satu huruf harus “dikorbankan”.

2. I dan J dianggap mirip secara historis

Secara historis, dalam alfabet Latin lama, I dan J tidak dibedakan.  J dianggap hanya variasi dari I.  

Karena itu, Playfair menggabungkan I dan J. Huruf J dihilangkan. Semua J dalam plaintext diganti 

menjadi I

Contoh:  JALAN → IALAN

3. Dampaknya ke proses enkripsi/dekripsi

Saat enkripsi: setiap J diperlakukan sebagai I. Saat dekripsi: huruf I bisa berarti I atau J, tergantung 

konteks bahasa. Ini tidak mengurangi keamanan, karena Playfair sudah bersifat substitusi digraf

Ambiguitas kecil ini masih bisa diselesaikan lewat konteks

4. Apakah bisa buang huruf lain?

Bisa saja secara teori, tapi  I/J paling masuk akal secara linguistik

Sudah jadi standar konvensi Playfair. Memudahkan implementasi dan pembacaan hasil

70

Sumber: ChatGPT

\end{document}
